\documentclass[12pt,a4paper]{article}

% ─── Encodage & langue ───────────────────────────────────────────────────────
\usepackage[utf8]{inputenc}
\usepackage[T1]{fontenc}
\usepackage[french]{babel}

% ─── Mise en page ────────────────────────────────────────────────────────────
\usepackage[top=2.5cm, bottom=2.5cm, left=2.8cm, right=2.8cm]{geometry}
\usepackage{setspace}
\onehalfspacing

% ─── Typographie ─────────────────────────────────────────────────────────────
\usepackage{microtype}

% ─── Mathématiques ───────────────────────────────────────────────────────────
\usepackage{amsmath, amssymb, amsthm, mathtools}
\usepackage{bm}

% ─── Couleurs & graphiques ───────────────────────────────────────────────────
\usepackage[dvipsnames]{xcolor}
\usepackage{tikz}
\usetikzlibrary{shapes.geometric, arrows.meta, positioning, fit, backgrounds}
\usepackage{pgfplots}
\pgfplotsset{compat=1.18}

% ─── Tableaux ────────────────────────────────────────────────────────────────
\usepackage{booktabs, array, tabularx, longtable, multirow}
\usepackage{colortbl}

% ─── Code source ─────────────────────────────────────────────────────────────
\usepackage{listings}

% ─── Liens & métadonnées ─────────────────────────────────────────────────────
\usepackage[colorlinks=true, linkcolor=NavyBlue, citecolor=ForestGreen,
            urlcolor=RoyalBlue, bookmarks=true]{hyperref}

% ─── En-têtes & pieds de page ────────────────────────────────────────────────
\usepackage{fancyhdr}
\usepackage{lastpage}

% ─── Environnements théorèmes ────────────────────────────────────────────────
\usepackage{thmtools}
\usepackage[framemethod=TikZ]{mdframed}

% ─── Boîtes colorées ─────────────────────────────────────────────────────────
\usepackage{tcolorbox}
\tcbuselibrary{skins, breakable, theorems}

% ─── Algorithmes ─────────────────────────────────────────────────────────────
\usepackage[french, ruled, linesnumbered]{algorithm2e}

% ─── Divers ──────────────────────────────────────────────────────────────────
\usepackage{enumitem}
\usepackage{caption}
\usepackage{subcaption}
\usepackage{float}
\usepackage{graphicx}

% ══════════════════════════════════════════════════════════════════════════════
%  COULEURS DU DOCUMENT
% ══════════════════════════════════════════════════════════════════════════════
\definecolor{AccentBlue}{RGB}{44, 91, 166}
\definecolor{AccentGreen}{RGB}{34, 139, 87}
\definecolor{LightBlue}{RGB}{230, 240, 255}
\definecolor{LightGreen}{RGB}{230, 248, 238}
\definecolor{LightGray}{RGB}{248, 249, 250}
\definecolor{DarkGray}{RGB}{52, 58, 64}
\definecolor{CodeBg}{RGB}{245, 247, 250}
\definecolor{TitleColor}{RGB}{15, 52, 96}

% ══════════════════════════════════════════════════════════════════════════════
%  STYLE DES LISTINGS
% ══════════════════════════════════════════════════════════════════════════════
\lstdefinestyle{pythonstyle}{
  language=Python,
  backgroundcolor=\color{CodeBg},
  basicstyle=\ttfamily\footnotesize,
  keywordstyle=\color{AccentBlue}\bfseries,
  stringstyle=\color{AccentGreen},
  commentstyle=\color{gray}\itshape,
  numberstyle=\tiny\color{gray},
  numbers=left,
  stepnumber=1,
  numbersep=8pt,
  frame=single,
  framerule=0.4pt,
  rulecolor=\color{AccentBlue!40},
  tabsize=4,
  showstringspaces=false,
  breaklines=true,
  captionpos=b,
  xleftmargin=15pt,
}
\lstset{style=pythonstyle}

% ══════════════════════════════════════════════════════════════════════════════
%  ENVIRONNEMENTS THÉORÈMES
% ══════════════════════════════════════════════════════════════════════════════
\tcbset{
  defstyle/.style={
    enhanced, breakable,
    colback=LightBlue, colframe=AccentBlue,
    fonttitle=\bfseries\small, title={#1},
    left=4pt, right=4pt, top=4pt, bottom=4pt,
  },
  propstyle/.style={
    enhanced, breakable,
    colback=LightGreen, colframe=AccentGreen,
    fonttitle=\bfseries\small, title={#1},
    left=4pt, right=4pt, top=4pt, bottom=4pt,
  },
  remarkstyle/.style={
    enhanced, breakable,
    colback=LightGray, colframe=DarkGray!60,
    fonttitle=\bfseries\small, title={#1},
    left=4pt, right=4pt, top=4pt, bottom=4pt,
  },
}

\newcounter{defcount}[section]
\newcounter{propcount}[section]

\newenvironment{definition}[1]{%
  \stepcounter{defcount}%
  \begin{tcolorbox}[defstyle={Définition \thesection.\thedefcount\ — #1}]
}{%
  \end{tcolorbox}
}

\newenvironment{proposition}[1]{%
  \stepcounter{propcount}%
  \begin{tcolorbox}[propstyle={Proposition \thesection.\thepropcount\ — #1}]
}{%
  \end{tcolorbox}
}

\newenvironment{remarque}[1]{%
  \begin{tcolorbox}[remarkstyle={Remarque — #1}]
}{%
  \end{tcolorbox}
}

\newenvironment{preuve}{%
  \noindent\textit{Démonstration.}\quad
}{%
  \hfill$\square$\par\medskip
}

% ══════════════════════════════════════════════════════════════════════════════
%  EN-TÊTES / PIEDS DE PAGE
% ══════════════════════════════════════════════════════════════════════════════
\pagestyle{fancy}
\fancyhf{}
\fancyhead[L]{\small\textcolor{AccentBlue}{\textbf{Feature Store Platform}}}
\fancyhead[R]{\small\textcolor{DarkGray}{LOMOFOUET NDONGMO Handrey Jauress}}
\fancyfoot[C]{\small\thepage\ / \pageref{LastPage}}
\renewcommand{\headrulewidth}{0.5pt}
\renewcommand{\footrulewidth}{0.3pt}

% ══════════════════════════════════════════════════════════════════════════════
%  DOCUMENT
% ══════════════════════════════════════════════════════════════════════════════
\begin{document}

% ─── Page de titre ───────────────────────────────────────────────────────────
\begin{titlepage}
\begin{tikzpicture}[remember picture, overlay]
  % Bandeau supérieur
  \fill[TitleColor] (current page.north west) rectangle
    ([yshift=-4cm]current page.north east);
  % Bandeau inférieur
  \fill[AccentBlue!80] (current page.south west) rectangle
    ([yshift=1.5cm]current page.south east);
  % Filet décoratif
  \fill[AccentGreen] ([yshift=-4cm]current page.north west) rectangle
    ([yshift=-4.15cm]current page.north east);
\end{tikzpicture}

\vspace*{2.5cm}
% \begin{center}
%   {\color{white}\Large\bfseries African Institute for Mathematical Sciences}\\[0.3cm]
%   {\color{white}\large AIMS Cameroun --- Dossier de Candidature}
% \end{center}

\vspace{3cm}
\begin{center}
  {\color{TitleColor}\rule{0.85\textwidth}{1.5pt}}\\[0.6cm]
  {\color{TitleColor}\Huge\bfseries Feature Store Platform}\\[0.4cm]
  {\color{AccentBlue}\LARGE Pipeline ML bout en bout :}\\[0.2cm]
  {\color{AccentBlue}\large Nettoyage, Feature Engineering,}\\[0.1cm]
  {\color{AccentBlue}\large Stockage versionné et API REST}\\[0.6cm]
  {\color{TitleColor}\rule{0.85\textwidth}{1.5pt}}
\end{center}

\vspace{1.5cm}
\begin{center}
  \begin{tabular}{ll}
    \textbf{Auteur}       & LOMOFOUET NDONGMO Handrey Jauress \\[4pt]
    \textbf{Domaines}     & Machine Learning, Statistiques, Génie Logiciel \\[4pt]
    \textbf{Technologies} & Python, Tkinter, SQLite, FastAPI, Pandas, Scikit-learn \\[4pt]
    \textbf{Document}     & Présentation technique et scientifique \\
    \textbf{URL Github} & \url{https://github.com/regent-ndongmo/feature-store-platform}} \\
  \end{tabular}
\end{center}

\vfill
\begin{center}
  {\color{white}\small Ce rapport présente la modélisation mathématique, les méthodes
  d'apprentissage automatique et l'ingénierie logicielle sous-jacentes à la
  \textbf{Feature Store Platform}, un outil complet de prétraitement de données pour
  l'apprentissage automatique.}
\end{center}
\vspace*{0.5cm}
\end{titlepage}

% ─── Table des matières ──────────────────────────────────────────────────────
\tableofcontents
\newpage

% ══════════════════════════════════════════════════════════════════════════════
\section{Introduction et Problématique}
% ══════════════════════════════════════════════════════════════════════════════

\subsection{Contexte}

Dans tout projet d'apprentissage automatique (\textit{machine learning}), la qualité
des données constitue le facteur déterminant de la performance des modèles.
La règle empirique, largement reconnue dans la communauté ML, stipule que 80\,\% du
temps d'un data scientist est consacré à la préparation des données
(\textit{data wrangling}), contre seulement 20\,\% à la modélisation proprement dite.

Ce constat soulève plusieurs défis fondamentaux :

\begin{itemize}[leftmargin=1.5em]
  \item \textbf{Hétérogénéité des données brutes :} valeurs manquantes, outliers,
    encodages incohérents, distributions asymétriques ;
  \item \textbf{Reproductibilité :} la même transformation appliquée par deux
    ingénieurs peut produire des résultats différents faute de standardisation ;
  \item \textbf{Fragmentation :} les \textit{features} calculées sont rarement
    centralisées, entraînant des recalculs redondants coûteux ;
  \item \textbf{Traçabilité :} l'absence de versioning des features rend difficile
    la comparaison d'expériences ML successives.
\end{itemize}

\subsection{Problème formel}

\begin{definition}{Pipeline de préparation de données}
Soit $\mathcal{D} = \{(\mathbf{x}_i, y_i)\}_{i=1}^{n}$ un jeu de données brut où
$\mathbf{x}_i \in \mathbb{R}^p$ est un vecteur de $p$ caractéristiques brutes et
$y_i \in \mathcal{Y}$ est l'étiquette cible. Le problème du prétraitement consiste
à trouver une transformation $\phi : \mathbb{R}^p \to \mathbb{R}^q$ telle que :
\[
  \phi(\mathbf{x}_i) = \mathbf{z}_i \in \mathbb{R}^q, \quad
  \text{avec } \mathbb{E}\bigl[\mathcal{L}(f(\mathbf{z}), y)\bigr]
  < \mathbb{E}\bigl[\mathcal{L}(f(\mathbf{x}), y)\bigr]
\]
où $\mathcal{L}$ est une fonction de perte, $f$ un modèle ML et $q \geq p$ ou
$q \leq p$ selon la nature des transformations appliquées.
\end{definition}

\subsection{Solution proposée}

La \textbf{Feature Store Platform} est un système complet qui résout ce problème en
articulant cinq composantes majeures :

\begin{enumerate}[leftmargin=1.5em]
  \item Un \textbf{moteur de nettoyage} basé sur des règles statistiques formelles
    (imputation, détection d'outliers, normalisation) ;
  \item Un \textbf{moteur de feature engineering} implémentant 11 transformations
    mathématiques fondamentales ;
  \item Un \textbf{feature store} avec stockage versionné (SQLite + CSV) et indexation ;
  \item Une \textbf{API REST} locale permettant la récupération standardisée des
    features pour l'entraînement de modèles ;
  \item Une \textbf{interface graphique Tkinter} guidant l'utilisateur à travers
    chaque étape du pipeline avec visualisations et pédagogie intégrée.
\end{enumerate}

% ══════════════════════════════════════════════════════════════════════════════
\section{Fondements Mathématiques du Nettoyage de Données}
% ══════════════════════════════════════════════════════════════════════════════

\subsection{Formalisation des données manquantes}

\begin{definition}{Mécanisme de données manquantes (Rubin, 1976)}
Soit $\mathbf{X} \in \mathbb{R}^{n \times p}$ la matrice de données complète et
$\mathbf{M} \in \{0,1\}^{n \times p}$ la matrice d'indicateurs de manquance
($M_{ij} = 1$ si $X_{ij}$ est observé). On distingue trois mécanismes :
\begin{itemize}
  \item \textbf{MCAR} (\textit{Missing Completely At Random}) :
    $P(\mathbf{M} \mid \mathbf{X}) = P(\mathbf{M})$ ;
  \item \textbf{MAR} (\textit{Missing At Random}) :
    $P(\mathbf{M} \mid \mathbf{X}) = P(\mathbf{M} \mid \mathbf{X}_{\text{obs}})$ ;
  \item \textbf{MNAR} (\textit{Missing Not At Random}) : dépendance aux valeurs
    manquantes elles-mêmes.
\end{itemize}
\end{definition}

\subsubsection{Imputation par la moyenne}

Pour une colonne $\mathbf{c} = (c_1, c_2, \ldots, c_n)^{\top}$ avec
$n_{\text{obs}}$ valeurs observées, l'estimateur de la moyenne est :
\begin{equation}
  \hat{\mu} = \bar{c} = \frac{1}{n_{\text{obs}}} \sum_{i \,:\, c_i \text{ observé}} c_i
\end{equation}

L'imputation substitue chaque valeur manquante par $\hat{\mu}$ :
\begin{equation}
  c_i^* = \begin{cases} c_i & \text{si } c_i \text{ est observé} \\
  \hat{\mu} & \text{sinon} \end{cases}
\end{equation}

\begin{remarque}{Biais de l'imputation par la moyenne}
Cette méthode est optimale au sens des moindres carrés sous hypothèse MCAR.
Elle minimise l'erreur quadratique moyenne $\text{MSE}(\hat{c}) =
\mathbb{E}[(c_i - \hat{\mu})^2]$. En revanche, elle réduit artificiellement
la variance empirique de la variable.
\end{remarque}

\subsubsection{Imputation par la médiane}

La médiane $\hat{m}$ est définie comme la solution du problème d'optimisation :
\begin{equation}
  \hat{m} = \arg\min_{\theta \in \mathbb{R}} \sum_{i \,:\, c_i \text{ observé}} |c_i - \theta|
\end{equation}

\begin{proposition}{Robustesse de la médiane}
Pour une distribution $F$ avec densité $f$, la variance asymptotique de la médiane
empirique est :
\[
  \text{Var}(\hat{m}) = \frac{1}{4 n f(\hat{m})^2}
\]
Contrairement à la moyenne, la médiane possède un point de rupture de $50\,\%$ :
elle reste un estimateur cohérent même si jusqu'à 50\,\% des observations sont
des outliers.
\end{proposition}

\subsubsection{Imputation par le mode (variables catégorielles)}

Pour une variable catégorielle $\mathbf{c}$ à valeurs dans $\mathcal{K}$, le mode
est l'estimateur du maximum de vraisemblance de la classe dominante :
\begin{equation}
  \hat{k} = \arg\max_{k \in \mathcal{K}} \frac{|\{i : c_i = k\}|}{n_{\text{obs}}}
  = \arg\max_{k \in \mathcal{K}} \hat{P}(C = k)
\end{equation}

\subsection{Détection et traitement des valeurs aberrantes}

\subsubsection{Méthode de l'écart interquartile (IQR)}

\begin{definition}{Bornes de Tukey}
Soit $Q_1$ et $Q_3$ les premier et troisième quartiles d'une distribution.
L'écart interquartile est $\text{IQR} = Q_3 - Q_1$. Les \textbf{bornes de Tukey}
définissent la région d'acceptation :
\begin{align}
  b_{\text{inf}} &= Q_1 - 1{,}5 \times \text{IQR} \\
  b_{\text{sup}} &= Q_3 + 1{,}5 \times \text{IQR}
\end{align}
Toute observation $c_i \notin [b_{\text{inf}}, b_{\text{sup}}]$ est classifiée comme
outlier.
\end{definition}

L'opération de \textit{clipping} remplace les outliers par les bornes :
\begin{equation}
  c_i^* = \max\!\bigl(b_{\text{inf}},\, \min(c_i,\, b_{\text{sup}})\bigr)
  = \text{clip}(c_i,\, b_{\text{inf}},\, b_{\text{sup}})
\end{equation}

\begin{remarque}{Justification statistique du facteur 1,5}
Sous hypothèse gaussienne $\mathcal{N}(\mu, \sigma^2)$, les bornes de Tukey
correspondent approximativement à $\mu \pm 2{,}7\sigma$, capturant 99{,}3\,\% de
la distribution. La probabilité qu'une observation gaussienne soit
déclarée outlier est inférieure à 0,7\,\%.
\end{remarque}

\subsubsection{Méthode du Z-score}

Le Z-score d'une observation est sa distance normalisée à la moyenne :
\begin{equation}
  z_i = \frac{c_i - \hat{\mu}}{\hat{\sigma}}, \quad
  \hat{\sigma} = \sqrt{\frac{1}{n-1}\sum_{j=1}^{n}(c_j - \hat{\mu})^2}
\end{equation}

Une observation est outlier si $|z_i| > z_{\alpha/2}$, où $z_{\alpha/2}$ est le
quantile de la loi normale standard au niveau $\alpha$ (typiquement $z_{0.025} = 1{,}96$
pour $\alpha = 5\,\%$, ou seuil $= 3$ dans notre implémentation).

\subsection{Encodage des variables catégorielles}

\subsubsection{Encodage par étiquettes (\textit{Label Encoding})}

Une variable catégorielle $C$ à $K$ modalités $\{k_1, k_2, \ldots, k_K\}$ est
transformée par une bijection entière :
\begin{equation}
  \text{LabelEnc} : \{k_1, \ldots, k_K\} \to \{0, 1, \ldots, K-1\}
\end{equation}

Cette représentation introduit un ordre artificiel entre les modalités, ce qui peut
biaiser les modèles fondés sur la distance (k-NN, SVM).

\subsubsection{Encodage \textit{one-hot}}

Pour éviter l'ordre artificiel, l'encodage one-hot représente chaque modalité par
un vecteur de la base canonique de $\mathbb{R}^K$ :
\begin{equation}
  \text{OneHot}(k_j) = \mathbf{e}_j \in \{0,1\}^K, \quad
  \text{tel que } (\mathbf{e}_j)_\ell = \mathbb{1}[\ell = j]
\end{equation}

La matrice d'encodage complète $\mathbf{E} \in \{0,1\}^{n \times K}$ vérifie
$\mathbf{E}\,\mathbf{1}_K = \mathbf{1}_n$ (chaque ligne somme à 1).

% ══════════════════════════════════════════════════════════════════════════════
\section{Mathématiques du Feature Engineering}
% ══════════════════════════════════════════════════════════════════════════════

Le feature engineering consiste à construire de nouvelles représentations
$\phi : \mathbb{R}^p \to \mathbb{R}^q$ maximisant l'information pertinente pour
la tâche d'apprentissage. Nous présentons les 11 transformations implémentées.

\subsection{Transformations de normalisation}

\subsubsection{Normalisation Min-Max (\textit{Min-Max Scaling})}

\begin{definition}{Normalisation Min-Max}
Pour une colonne $\mathbf{c} \in \mathbb{R}^n$, la normalisation min-max
projette les valeurs dans l'intervalle $[0, 1]$ :
\begin{equation}
  \tilde{c}_i = \frac{c_i - c_{\min}}{c_{\max} - c_{\min}}, \quad
  c_{\min} = \min_j c_j, \quad c_{\max} = \max_j c_j
\end{equation}
\end{definition}

\begin{proposition}{Conservation des rapports}
La transformation min-max est une homothétie affine. Elle préserve l'ordre,
les rapports de distances :
\[
  \frac{\tilde{c}_i - \tilde{c}_j}{\tilde{c}_k - \tilde{c}_l}
  = \frac{c_i - c_j}{c_k - c_l}, \qquad \forall\, i,j,k,l
\]
et la structure géométrique de l'espace des données.
\end{proposition}

\subsubsection{Standardisation Z-score (\textit{Standard Scaling})}

\begin{definition}{Standardisation}
La standardisation centre et réduit une variable :
\begin{equation}
  \tilde{c}_i = \frac{c_i - \hat{\mu}}{\hat{\sigma}},
  \quad \hat{\mu} = \frac{1}{n}\sum_i c_i,
  \quad \hat{\sigma}^2 = \frac{1}{n-1}\sum_i (c_i - \hat{\mu})^2
\end{equation}
\end{definition}

La variable standardisée vérifie $\hat{\mathbb{E}}[\tilde{c}] = 0$ et
$\hat{\text{Var}}[\tilde{c}] = 1$. Cette propriété est essentielle pour les
algorithmes sensibles à l'échelle (régression Ridge, PCA, SVM avec noyau RBF).

\subsection{Transformations de stabilisation de variance}

\subsubsection{Transformation logarithmique}

\begin{definition}{Transformation log (Box-Cox, $\lambda=0$)}
Pour $c_i > 0$, la transformation log est définie par :
\begin{equation}
  \tilde{c}_i = \log(1 + c_i)
\end{equation}
L'ajout de la constante $+1$ assure la définition en $c_i = 0$.
\end{definition}

\begin{proposition}{Stabilisation de variance}
Si $c_i \sim \text{LogNormal}(\mu, \sigma^2)$, alors
$\log(c_i) \sim \mathcal{N}(\mu, \sigma^2)$.
La transformation log est l'outil de choix pour les variables à queue droite
épaisse (revenus, prix, populations), car elle rend leur distribution
approximativement gaussienne, condition souvent requise par les modèles
paramétriques.
\end{proposition}

La dérivée $\frac{d}{dc}\log(1+c) = \frac{1}{1+c}$ décroît avec $c$ :
la transformation compresse les grandes valeurs plus que les petites,
réduisant ainsi l'asymétrie (\textit{skewness}) positive.

\subsubsection{Transformation racine carrée}

Cas particulier de la famille de transformations puissance
$\phi_\lambda(c) = c^\lambda$ avec $\lambda = 1/2$ :
\begin{equation}
  \tilde{c}_i = \sqrt{c_i}
\end{equation}

Cette transformation est appropriée pour les variables de comptage
(distributions de Poisson) où $\text{Var}(c) \propto \mathbb{E}[c]$.
La transformation stabilise la variance : $\text{Var}(\sqrt{c}) \approx \frac{1}{4}$.

\subsection{Discrétisation et binarisation}

\subsubsection{Binarisation par seuillage}

\begin{definition}{Binarisation}
Pour un seuil $\theta \in \mathbb{R}$, la binarisation est la fonction indicatrice :
\begin{equation}
  \tilde{c}_i = \mathbb{1}[c_i \geq \theta] =
  \begin{cases} 1 & \text{si } c_i \geq \theta \\ 0 & \text{sinon} \end{cases}
\end{equation}
\end{definition}

Cette transformation permet de créer des variables booléennes à partir de variables
continues (par exemple : kilométrage élevé / faible, prix premium / standard).
Le seuil optimal $\theta^*$ peut être déterminé par maximisation de l'information
mutuelle $I(C; Y)$ entre la variable binarisée et la cible :
\begin{equation}
  \theta^* = \arg\max_\theta I(\mathbb{1}[C \geq \theta];\, Y)
\end{equation}

\subsubsection{Discrétisation par intervalles égaux (\textit{Binning})}

La discrétisation partitionne l'espace des valeurs en $K$ intervalles :
\begin{equation}
  \mathcal{I}_k = \left[c_{\min} + (k-1)\cdot\frac{c_{\max}-c_{\min}}{K},\;
  c_{\min} + k\cdot\frac{c_{\max}-c_{\min}}{K}\right), \quad k = 1,\ldots,K
\end{equation}

Chaque valeur est remplacée par l'indice de son intervalle :
\begin{equation}
  \tilde{c}_i = \sum_{k=1}^{K} k \cdot \mathbb{1}[c_i \in \mathcal{I}_k]
\end{equation}

\subsection{Interactions et combinaisons de features}

\subsubsection{Ratio entre deux variables}

\begin{definition}{Feature de ratio}
Pour deux colonnes $\mathbf{a}, \mathbf{b} \in \mathbb{R}^n$ avec
$b_i \neq 0$, le ratio est :
\begin{equation}
  r_i = \frac{a_i}{b_i + \varepsilon}, \quad \varepsilon > 0 \text{ terme de régularisation}
\end{equation}
\end{definition}

Cette transformation est fondamentale en finance (ratio prix/bénéfice),
en physique (densité = masse/volume) et en ML pour créer des features
invariantes aux changements d'échelle.

\subsubsection{Interaction polynomiale}

Le produit de deux features capture les interactions non linéaires :
\begin{equation}
  p_i = a_i \times b_i
\end{equation}

Dans le cadre de la régression linéaire généralisée, cette feature correspond
au terme d'interaction dans le développement polynomial de degré 2 :
\begin{equation}
  f(\mathbf{x}) = \beta_0 + \beta_1 a + \beta_2 b + \beta_{12} ab + \cdots
\end{equation}

\subsubsection{Différence entre deux variables}

\begin{equation}
  d_i = a_i - b_i
\end{equation}

Cette transformation permet de capturer des écarts significatifs (par exemple :
différence entre le prix demandé et le prix de marché estimé).

\subsection{Pipeline de feature engineering}

L'enchaînement complet des transformations peut être formalisé comme une composition
de fonctions :

\begin{equation}
  \phi = \phi_K \circ \phi_{K-1} \circ \cdots \circ \phi_1
\end{equation}

où chaque $\phi_k$ est l'une des transformations définies ci-dessus.
Le pipeline maintient une file d'opérations ordonnées garantissant la
reproductibilité.

% ══════════════════════════════════════════════════════════════════════════════
\section{Le Feature Store : Modélisation et Fondements}
% ══════════════════════════════════════════════════════════════════════════════

\subsection{Définition formelle}

\begin{definition}{Feature Store}
Un \textbf{feature store} est un système de stockage centralisé $\mathcal{F}$ défini
par un triplet :
\begin{equation}
  \mathcal{F} = (\mathcal{G}, \mathcal{V}, \mathcal{I})
\end{equation}
où :
\begin{itemize}
  \item $\mathcal{G}$ est l'ensemble des \textbf{groupes de features},
    $\mathcal{G} = \{G_1, G_2, \ldots, G_m\}$ ;
  \item $\mathcal{V}(G_k) = \{v_1^k, v_2^k, \ldots\}$ est l'ensemble des
    \textbf{versions} du groupe $G_k$ ;
  \item $\mathcal{I}$ est une fonction d'indexation reliant entités et features :
    $\mathcal{I} : \text{entity\_id} \times \text{group} \times \text{version} \to \mathbb{R}^q$.
\end{itemize}
\end{definition}

\subsection{Versioning et gestion de l'historique}

Le versioning dans notre implémentation suit une logique de séquence monotone
croissante. Lors de l'enregistrement d'un groupe $G_k$ :

\begin{equation}
  v_{\text{new}}(G_k) = \max\bigl(\{v : v \in \mathcal{V}(G_k)\}\bigr) + 1
\end{equation}

Cette stratégie garantit l'immuabilité des versions passées
(\textit{append-only semantics}), permettant la comparaison d'expériences ML
à différentes étapes du prétraitement.

\subsection{Récupération par filtrage (Point-in-Time Lookup)}

\begin{definition}{Requête de récupération}
Une requête de récupération est définie par un triplet
$q = (\text{group}, F, E)$ où $F \subseteq \{f_1,\ldots,f_q\}$ est
l'ensemble des features demandées et $E \subseteq \{\text{entity\_ids}\}$
est l'ensemble des entités. Le résultat est :
\begin{equation}
  \text{fetch}(q) = \pi_F\!\bigl(\sigma_E(D_{\text{group}})\bigr)
\end{equation}
où $\pi_F$ est la projection sur les colonnes $F$ et $\sigma_E$ est la sélection
des lignes d'entités $E$, au sens de l'algèbre relationnelle.
\end{definition}

\subsection{Modèle relationnel de persistance}

Les données sont persistées dans deux couches complémentaires :

\begin{center}
\begin{tikzpicture}[
  every node/.style={font=\small},
  table/.style={draw, rounded corners=3pt, fill=LightBlue, minimum width=5cm},
  arrow/.style={-Stealth, thick, AccentBlue}
]
  \node[table, minimum height=2.8cm] (meta) at (0,0) {
    \begin{tabular}{l}
      \textbf{feature\_groups (SQLite)}\\
      \hline
      id : INTEGER PK\\
      name : TEXT\\
      description : TEXT\\
      entity\_col : TEXT\\
      feature\_cols : TEXT (JSON)\\
      version : INTEGER\\
      csv\_path : TEXT\\
      created\_at : TEXT\\
    \end{tabular}
  };
  \node[table, minimum height=1.5cm] (csv) at (7,0) {
    \begin{tabular}{l}
      \textbf{<group>\_v<n>.csv}\\
      \hline
      entity\_col | feature$_1$ | \ldots | feature$_q$\\
      \ldots
    \end{tabular}
  };
  \draw[arrow] (meta.east) -- node[above]{\footnotesize csv\_path} (csv.west);
\end{tikzpicture}
\end{center}

\subsection{Complexité des opérations}

Soit $n$ le nombre de lignes, $p$ le nombre de colonnes totales,
$q \leq p$ le nombre de features sélectionnées et $|E|$ la taille de l'ensemble
d'entités requêtées.

\begin{proposition}{Complexité des opérations du Feature Store}
\begin{align}
  \text{Écriture (save)} &: \mathcal{O}(n \cdot p) \\
  \text{Lecture complète (fetch all)} &: \mathcal{O}(n \cdot q) \\
  \text{Lecture filtrée (fetch subset)} &: \mathcal{O}(|E| \cdot q) \\
  \text{Listage des groupes} &: \mathcal{O}(m)
\end{align}
où $m = |\mathcal{G}|$ est le nombre de groupes enregistrés.
\end{proposition}

% ══════════════════════════════════════════════════════════════════════════════
\section{Statistiques Descriptives et Analyse Exploratoire}
% ══════════════════════════════════════════════════════════════════════════════

\subsection{Mesures de tendance centrale et de dispersion}

Le module de chargement calcule automatiquement un profil statistique de chaque
colonne numérique.

\begin{definition}{Statistiques d'ordre}
Pour un échantillon $\mathbf{c} = (c_{(1)}, c_{(2)}, \ldots, c_{(n)})^{\top}$
trié par ordre croissant, les statistiques d'ordre sont :
\begin{itemize}
  \item \textbf{Minimum} : $c_{(1)}$ ;
  \item \textbf{Maximum} : $c_{(n)}$ ;
  \item \textbf{Médiane} : $c_{(n/2)}$ si $n$ pair,
    $\frac{1}{2}(c_{(\lfloor n/2 \rfloor)} + c_{(\lceil n/2 \rceil)})$ sinon ;
  \item \textbf{Quantile $\alpha$} :
    $Q_\alpha = c_{(\lfloor \alpha n \rfloor + 1)}$.
\end{itemize}
\end{definition}

\subsection{Coefficient d'asymétrie (\textit{Skewness})}

\begin{equation}
  \gamma_1 = \frac{\hat{\mu}_3}{\hat{\sigma}^3},
  \quad \hat{\mu}_3 = \frac{1}{n}\sum_{i=1}^n (c_i - \bar{c})^3
\end{equation}

\begin{itemize}
  \item $\gamma_1 > 0$ : distribution asymétrique à droite (queue à droite) ;
  \item $\gamma_1 < 0$ : distribution asymétrique à gauche ;
  \item $\gamma_1 \approx 0$ : distribution symétrique.
\end{itemize}

Le coefficient de skewness guide le choix de la transformation appropriée :
si $|\gamma_1| > 1$, la transformation logarithmique ou racine carrée est recommandée.

\subsection{Taux de valeurs manquantes}

\begin{equation}
  \text{NaN\_rate}(\mathbf{c}) = \frac{|\{i : c_i \text{ est manquant}\}|}{n}
  = 1 - \frac{n_{\text{obs}}}{n}
\end{equation}

Cette métrique guide la stratégie d'imputation :
$\text{NaN\_rate} < 5\,\%$ : imputation acceptable ;
$5\,\% \leq \text{NaN\_rate} < 30\,\%$ : imputation avec prudence ;
$\text{NaN\_rate} \geq 30\,\%$ : suppression de la colonne recommandée.

% ══════════════════════════════════════════════════════════════════════════════
\section{Architecture Logicielle}
% ══════════════════════════════════════════════════════════════════════════════

\subsection{Principes de génie logiciel appliqués}

L'architecture respecte les principes fondamentaux du génie logiciel moderne :

\subsubsection{Séparation des préoccupations (Separation of Concerns)}

Le système est découpé en couches indépendantes communiquant via des interfaces
bien définies :

\begin{center}
\begin{tikzpicture}[
  layer/.style={draw, rounded corners=4pt, minimum width=10cm,
                minimum height=0.9cm, font=\small\bfseries},
  arrow/.style={-Stealth, thick, DarkGray}
]
  \node[layer, fill=AccentBlue!20, draw=AccentBlue] (ui)
    at (0, 4.5) {Couche Présentation --- Interface Tkinter + Matplotlib};
  \node[layer, fill=AccentGreen!20, draw=AccentGreen] (api)
    at (0, 3.2) {Couche API --- FastAPI / HTTP (port 5050)};
  \node[layer, fill=orange!20, draw=orange!80] (biz)
    at (0, 1.9) {Couche Métier --- DataCleaner, FeatureEngineer};
  \node[layer, fill=purple!15, draw=purple!60] (data)
    at (0, 0.6) {Couche Données --- FeatureStore (SQLite + CSV)};

  \draw[arrow] (ui) -- (api)
    node[midway, right=0.3cm]{\footnotesize HTTP/REST};
  \draw[arrow] (api) -- (biz)
    node[midway, right=0.3cm]{\footnotesize appels de services};
  \draw[arrow] (biz) -- (data)
    node[midway, right=0.3cm]{\footnotesize requêtes persistance};
\end{tikzpicture}
\end{center}

\subsubsection{Principes SOLID}

\begin{itemize}[leftmargin=1.5em]
  \item \textbf{S} (\textit{Single Responsibility}) : \texttt{DataCleaner} ne
    connaît pas le store ; \texttt{FeatureStore} ne connaît pas les transformations ;
  \item \textbf{O} (\textit{Open/Closed}) : de nouvelles règles de nettoyage ou
    transformations s'ajoutent sans modifier le code existant ;
  \item \textbf{D} (\textit{Dependency Inversion}) : l'interface graphique
    dépend des contrats de service, non des implémentations concrètes.
\end{itemize}

\subsection{Composants principaux}

\subsubsection{Le moteur de nettoyage (\texttt{cleaner.py})}

\begin{lstlisting}[caption={Interface du moteur de nettoyage},
                   label={lst:cleaner}]
class DataCleaner:
    """
    Moteur de nettoyage : 11 regles statistiques formelles.
    Chaque operation est une transformation pure :
    DataFrame -> DataFrame (pas d'effets de bord).
    """

    OPERATIONS = {
        'fill_mean'      : 'Imputation par la moyenne',
        'fill_median'    : 'Imputation par la mediane',
        'fill_mode'      : 'Imputation par le mode',
        'fill_constant'  : 'Imputation par constante',
        'drop_na_rows'   : 'Suppression des lignes avec NaN',
        'drop_column'    : 'Suppression de colonne',
        'clip_outliers'  : 'Ecrêtage outliers (IQR)',
        'label_encode'   : 'Encodage par etiquettes',
        'one_hot_encode' : 'Encodage one-hot',
        'to_datetime'    : 'Conversion en datetime',
        'strip_spaces'   : 'Nettoyage des espaces',
    }

    def apply_rule(self, df: pd.DataFrame,
                   action: str, col: str,
                   value=None) -> pd.DataFrame:
        """Applique une regle de nettoyage."""
        handler = getattr(self, f'_rule_{action}')
        return handler(df.copy(), col, value)

    def _rule_clip_outliers(self, df, col, _):
        Q1 = df[col].quantile(0.25)
        Q3 = df[col].quantile(0.75)
        IQR = Q3 - Q1
        lb, ub = Q1 - 1.5*IQR, Q3 + 1.5*IQR
        df[col] = df[col].clip(lb, ub)   # c* = clip(c, lb, ub)
        return df
\end{lstlisting}

\subsubsection{Le moteur de feature engineering (\texttt{engineer.py})}

\begin{lstlisting}[caption={Implémentation des transformations clés},
                   label={lst:engineer}]
class FeatureEngineer:
    """
    11 transformations de feature engineering.
    Chaque transformation retourne un DataFrame enrichi
    avec la nouvelle colonne calculee.
    """

    def log_transform(self, df, col, new_col):
        # tilde_c_i = log(1 + c_i)
        df[new_col] = np.log1p(df[col])
        return df

    def normalize(self, df, col, new_col):
        # tilde_c_i = (c_i - c_min) / (c_max - c_min)
        mn, mx = df[col].min(), df[col].max()
        df[new_col] = (df[col] - mn) / (mx - mn + 1e-9)
        return df

    def standardize(self, df, col, new_col):
        # tilde_c_i = (c_i - mu) / sigma
        mu, sigma = df[col].mean(), df[col].std()
        df[new_col] = (df[col] - mu) / (sigma + 1e-9)
        return df

    def binarize(self, df, col, new_col, threshold=0.0):
        # tilde_c_i = 1[c_i >= threshold]
        df[new_col] = (df[col] >= threshold).astype(int)
        return df
\end{lstlisting}

\subsubsection{Le Feature Store (\texttt{feature\_store.py})}

\begin{lstlisting}[caption={Gestion de la persistance et du versioning},
                   label={lst:store}]
class FeatureStore:
    """
    Stockage centralise des features : SQLite (metadonnees)
    + CSV (donnees). Gestion automatique du versioning.
    """

    def save_feature_group(self, name, df, entity_col,
                           feature_cols, description=""):
        # v_new = max(versions existantes) + 1
        existing = self._get_versions(name)
        version = max(existing, default=0) + 1

        csv_path = f"store/data/{name}_v{version}.csv"
        df[feature_cols].to_csv(csv_path, index=False)

        # Persistance metadonnees en SQLite
        self.conn.execute("""
            INSERT INTO feature_groups
            (name, description, entity_col, feature_cols,
             version, csv_path, created_at)
            VALUES (?,?,?,?,?,?,?)
        """, (name, description, entity_col,
              json.dumps(feature_cols), version,
              csv_path, datetime.now().isoformat()))
        return version

    def fetch(self, group_name, features=None,
              entity_ids=None):
        # Projection pi_F et selection sigma_E
        df = pd.read_csv(csv_path)
        if features:
            df = df[features]                   # pi_F
        if entity_ids:
            df = df[df[entity_col].isin(entity_ids)]  # sigma_E
        return df
\end{lstlisting}

\subsection{L'API REST (\texttt{server.py})}

L'API expose le feature store via un serveur HTTP minimal conforme au style
architectural REST (\textit{Representational State Transfer}).

\begin{center}
\renewcommand{\arraystretch}{1.4}
\begin{tabular}{llp{6cm}}
  \toprule
  \textbf{Méthode} & \textbf{Endpoint} & \textbf{Description} \\
  \midrule
  \texttt{GET}    & \texttt{/api/feature-groups}
                  & Lister tous les groupes enregistrés \\
  \texttt{GET}    & \texttt{/api/feature-groups/<nom>}
                  & Métadonnées d'un groupe \\
  \texttt{POST}   & \texttt{/api/feature-groups/<nom>/fetch}
                  & Récupérer les features (avec filtres) \\
  \texttt{DELETE} & \texttt{/api/feature-groups/<nom>}
                  & Supprimer un groupe \\
  \texttt{GET}    & \texttt{/health}
                  & État du serveur \\
  \bottomrule
\end{tabular}
\end{center}

Un exemple de réponse de l'endpoint \texttt{fetch} :

\begin{lstlisting}[language={},caption={Réponse JSON de /api/feature-groups/vehicles/fetch}]
{
  "group": "vehicles",
  "version": 2,
  "features": ["price", "odometer", "log_price"],
  "entity_count": 150,
  "data": [
    {"id": 0, "price": 12500, "odometer": 82000, "log_price": 9.43},
    {"id": 1, "price": 8900,  "odometer": 115000, "log_price": 9.09}
  ]
}
\end{lstlisting}

% ══════════════════════════════════════════════════════════════════════════════
\section{Visualisations et Pédagogie Intégrée}
% ══════════════════════════════════════════════════════════════════════════════

\subsection{Graphiques contextuels (Matplotlib + PIL)}

L'interface génère des graphiques en temps réel en utilisant le backend
Matplotlib \texttt{Agg} (rendu \textit{off-screen} en PNG, converti en
\texttt{PhotoImage} Tkinter via PIL).

Pour chaque opération de nettoyage, le graphique illustre l'effet statistique :

\begin{itemize}
  \item \textbf{Imputation (fill\_mean/median)} : histogramme de distribution
    avec ligne verticale indiquant $\hat{\mu}$ ou $\hat{m}$ ;
  \item \textbf{Outliers (clip)} : histogramme avec zones colorées distinguant
    les valeurs inliers (bleu) et outliers (rouge) définies par les bornes de Tukey ;
  \item \textbf{Suppression (drop\_na)} : diagramme en barres montrant
    $n_{\text{obs}}$ vs $n_{\text{NaN}}$ avec le ratio NaN.
\end{itemize}

Pour le feature engineering, chaque graphique est un diagramme \textbf{Avant ↔ Après}
(deux sous-graphiques côte à côte) illustrant la transformation de la distribution.

\subsection{Indicateur de progression du workflow}

L'interface implémente une barre de progression en 5 étapes
$\{①\text{Charger},\, ②\text{Nettoyer},\, ③\text{Transformer},\,
④\text{Stocker},\, ⑤\text{Utiliser}\}$ mise à jour automatiquement.

Cette progression peut être formalisée comme un automate fini déterministe
$\mathcal{A} = (Q, \Sigma, \delta, q_0, F)$ où :
\begin{itemize}
  \item $Q = \{0, 1, 2, 3, 4\}$ sont les états (étapes) ;
  \item $\Sigma = \{\text{load}, \text{clean}, \text{engineer}, \text{store},
    \text{use}\}$ les événements déclencheurs ;
  \item $\delta : Q \times \Sigma \to Q$ la fonction de transition ;
  \item $q_0 = 0$ l'état initial, $F = \{4\}$ l'état final.
\end{itemize}

% ══════════════════════════════════════════════════════════════════════════════
\section{Résultats et Évaluation Expérimentale}
% ══════════════════════════════════════════════════════════════════════════════

\subsection{Jeu de données de référence}

Le jeu de données de démonstration contient $n = 300$ entrées de véhicules d'occasion
avec les colonnes suivantes :

\begin{center}
\renewcommand{\arraystretch}{1.3}
\begin{tabular}{llcc}
  \toprule
  \textbf{Colonne} & \textbf{Type} & \textbf{NaN\,\%} & \textbf{Transformation appliquée} \\
  \midrule
  \texttt{price}       & numérique & 3.3\,\%  & log\_transform \\
  \texttt{odometer}    & numérique & 5.0\,\%  & normalize \\
  \texttt{year}        & numérique & 0.0\,\%  & binarize ($\theta=2018$) \\
  \texttt{condition}   & catég.    & 8.3\,\%  & label\_encode \\
  \texttt{fuel}        & catég.    & 1.7\,\%  & one\_hot\_encode \\
  \bottomrule
\end{tabular}
\end{center}

\subsection{Exemple numérique : transformation log}

Avant transformation, la variable \texttt{price} a une distribution asymétrique :
\[
  \hat{\mu}_{\text{price}} = 12\,847\,\$,\quad
  \hat{\sigma}_{\text{price}} = 9\,213\,\$,\quad
  \gamma_1 = 2{,}1 \quad \text{(forte asymétrie droite)}
\]

Après transformation log :
\[
  \hat{\mu}_{\log(1+\text{price})} = 9{,}27,\quad
  \hat{\sigma}_{\log(1+\text{price})} = 0{,}74,\quad
  \gamma_1 = 0{,}18 \quad \text{(quasi-symétrique)}
\]

La transformation a réduit le coefficient d'asymétrie d'un facteur
$2{,}1 / 0{,}18 \approx 11{,}7$.

\subsection{Exemple numérique : détection d'outliers IQR}

Pour \texttt{odometer} :
\[
  Q_1 = 52\,000,\quad Q_3 = 138\,000,\quad
  \text{IQR} = 86\,000
\]
\[
  b_{\text{inf}} = 52\,000 - 1{,}5 \times 86\,000 = -77\,000 \approx 0
\]
\[
  b_{\text{sup}} = 138\,000 + 1{,}5 \times 86\,000 = 267\,000
\]

Sur 300 véhicules, 8 ont un kilométrage $> 267\,000$\,km et sont écrêtés.

\subsection{Performance du pipeline complet}

Le pipeline complet sur 300 lignes et 41 tests unitaires :

\begin{center}
\renewcommand{\arraystretch}{1.3}
\begin{tabular}{lcc}
  \toprule
  \textbf{Composant} & \textbf{Tests} & \textbf{Taux de succès} \\
  \midrule
  DataCleaner    & 13 & 100\,\% \\
  FeatureEngineer & 11 & 100\,\% \\
  FeatureStore   & 11 & 100\,\% \\
  API REST       & 5  & 100\,\% \\
  End-to-End     & 1  & 100\,\% \\
  \midrule
  \textbf{Total} & \textbf{41} & \textbf{100\,\%} \\
  \bottomrule
\end{tabular}
\end{center}

% ══════════════════════════════════════════════════════════════════════════════
\section{Génie Logiciel --- Méthodologie et Qualité}
% ══════════════════════════════════════════════════════════════════════════════

\subsection{Structure du projet}

\begin{lstlisting}[language={}, caption={Arborescence du projet}]
feature_store_app/
|-- app.py                  # Interface Tkinter (1537 lignes)
|-- core/
|   |-- cleaner.py          # 11 regles de nettoyage
|   `-- engineer.py         # 11 transformations FE
|-- store/
|   |-- feature_store.py    # Backend SQLite + CSV
|   `-- data/               # *.db + *.csv (versions)
|-- api/
|   `-- server.py           # Serveur HTTP port 5050
|-- tests/
|   `-- run_tests.py        # 41 tests (pytest)
|-- requirements.txt
`-- README.md
\end{lstlisting}

\subsection{Algorithme du pipeline complet}

\begin{algorithm}[H]
\caption{Pipeline Feature Store --- Workflow complet}
\KwIn{Fichier de données $F$ (CSV / Excel / Parquet)}
\KwOut{Feature group versionnée $G_k^{(v)}$ disponible via API}

\BlankLine
\textbf{Étape 1 : Chargement}\\
Lire $F$ avec Pandas $\Rightarrow$ DataFrame $\mathbf{X}^{(0)} \in \mathbb{R}^{n \times p}$\\
Calculer le profil statistique : $(\hat{\mu}_j, \hat{\sigma}_j, \text{NaN\_rate}_j)$ pour $j=1,\ldots,p$\\

\BlankLine
\textbf{Étape 2 : Nettoyage}\\
\ForEach{règle de nettoyage $(a, c, v)$ dans la file $\mathcal{Q}$}{
  $\mathbf{X} \leftarrow \text{apply\_rule}(\mathbf{X}, a, c, v)$
}
$\mathbf{X}^{(1)} \leftarrow \mathbf{X}$ \quad \textit{(DataFrame nettoyé)}\\

\BlankLine
\textbf{Étape 3 : Feature Engineering}\\
\ForEach{transformation $(\phi_k, \text{col}, \text{params})$ dans la file}{
  $\mathbf{X} \leftarrow \phi_k(\mathbf{X})$ \quad \textit{(ajout colonne)}
}
$\mathbf{X}^{(2)} \leftarrow \mathbf{X}$ \quad \textit{(DataFrame enrichi, } $q \geq p$\textit{)}\\

\BlankLine
\textbf{Étape 4 : Stockage}\\
$v \leftarrow \max(\mathcal{V}(G_k)) + 1$\\
Persister $\pi_F(\mathbf{X}^{(2)})$ en CSV et métadonnées en SQLite\\

\BlankLine
\textbf{Étape 5 : Exposition API}\\
Répondre aux requêtes \texttt{POST /api/feature-groups/$G_k$/fetch}\\
Retourner $\sigma_E(\pi_F(\mathbf{X}^{(2)}))$ au format JSON
\end{algorithm}

\subsection{Processus de développement agile}

Le projet a été conduit selon une approche itérative :

\begin{enumerate}[leftmargin=1.5em]
  \item \textbf{Spécification} : définition des 11 règles de nettoyage et des
    11 transformations FE avec leurs fondements mathématiques ;
  \item \textbf{Conception} : architecture n-tiers, modèle de données SQLite,
    contrat d'API REST ;
  \item \textbf{Implémentation} :
    v1.0 (noyau ML + backend store + API) $\to$ v2.0 (interface Tkinter premium
    + graphiques contextuels + tooltips pédagogiques) ;
  \item \textbf{Validation} : 41 tests unitaires et d'intégration (pytest).
\end{enumerate}

\subsection{Qualité du code}

\begin{itemize}[leftmargin=1.5em]
  \item \textbf{Tests} : couverture à 100\,\% des 41 cas de test ;
  \item \textbf{Documentation} : docstrings pour toutes les classes et méthodes ;
  \item \textbf{Gestion des erreurs} : chaque transformation gère les cas
    limites (division par zéro, colonne non numérique, fichier absent) ;
  \item \textbf{Validation des données} : vérification du type de colonne avant
    application de toute transformation numérique ;
  \item \textbf{Thread-safety} : l'API REST et l'interface graphique tournent
    dans des threads séparés avec verrous mutex sur l'accès SQLite.
\end{itemize}

% ══════════════════════════════════════════════════════════════════════════════
\section{Limites et Perspectives}
% ══════════════════════════════════════════════════════════════════════════════

\begin{center}
\renewcommand{\arraystretch}{1.4}
\begin{tabular}{p{5.5cm}p{7cm}}
  \toprule
  \textbf{Limite actuelle} & \textbf{Extension envisagée} \\
  \midrule
  Scalabilité : chargement en mémoire & Traitement par blocs (\textit{chunked processing})
    avec Dask ou Spark \\
  Imputation univariée & Imputation multivariée : MICE, KNN-imputation
    ($\hat{c}_i = \frac{\sum_{j \in kNN} w_j c_j}{\sum w_j}$) \\
  Seuils d'outliers fixes (IQR) & Détection adaptative : Isolation Forest,
    DBSCAN, Autoencodeurs \\
  API synchrone monothread & API asynchrone avec \texttt{asyncio} + FastAPI \\
  Transformations indépendantes & Pipeline scikit-learn exportable
    (\texttt{Pipeline}, \texttt{ColumnTransformer}) \\
  Stockage local SQLite + CSV & Backends distribués : Redis Feature Store,
    Feast, Tecton \\
  \bottomrule
\end{tabular}
\end{center}

% ══════════════════════════════════════════════════════════════════════════════
\section{Conclusion}
% ══════════════════════════════════════════════════════════════════════════════

Ce projet démontre la synergie rigoureuse entre plusieurs disciplines
mathématiques et informatiques :

\begin{itemize}[leftmargin=1.5em]
  \item \textbf{Statistiques} : théorie de l'estimation (médiane, moyenne,
    quartiles), robustesse des estimateurs, asymétrie, mécanismes de données
    manquantes de Rubin, complexité algorithmique ;
  \item \textbf{Algèbre linéaire} : espaces vectoriels $\mathbb{R}^p$, transformations
    linéaires (standardisation), produits cartésiens (one-hot encoding),
    projections (feature selection) ;
  \item \textbf{Analyse mathématique} : transformations de stabilisation de variance
    (log, racine carrée), fonctions monotones, inégalités ;
  \item \textbf{Apprentissage automatique} : pipeline de prétraitement,
    feature engineering, algèbre relationnelle pour le feature store ;
  \item \textbf{Génie logiciel} : architecture n-tiers, principes SOLID,
    API REST, tests automatisés, versioning, thread-safety.
\end{itemize}

\begin{tcolorbox}[enhanced, colback=TitleColor!8, colframe=TitleColor,
  fonttitle=\bfseries, title={Portée scientifique --- Lien avec AIMS}]
Ce travail s'inscrit directement dans les thématiques d'AIMS : optimisation,
analyse matricielle, apprentissage statistique et algorithmique.
Il constitue une application concrète de l'algèbre linéaire et des statistiques
à des problèmes réels de préparation de données ML, ouvrant la voie à des
extensions mathématiques plus sophistiquées (optimisation convexe des paramètres
d'imputation, théorie de l'information pour la sélection de features,
apprentissage de métriques de distance).

Ce projet, entièrement conçu, architecturé et implémenté de manière autonome,
témoigne d'une maîtrise transversale allant de la modélisation mathématique
formelle à la livraison d'un logiciel professionnel complet et testé.
\end{tcolorbox}

\newpage
% ══════════════════════════════════════════════════════════════════════════════
%  RÉFÉRENCES
% ══════════════════════════════════════════════════════════════════════════════
\begin{thebibliography}{99}

\bibitem{rubin1976}
Rubin, D.~B. (1976).
\textit{Inference and Missing Data}.
Biometrika, 63(3), 581--592.

\bibitem{tukey1977}
Tukey, J.~W. (1977).
\textit{Exploratory Data Analysis}.
Addison-Wesley.

\bibitem{boxcox1964}
Box, G.~E.~P. \& Cox, D.~R. (1964).
An analysis of transformations.
\textit{Journal of the Royal Statistical Society, Series B}, 26(2), 211--252.

\bibitem{sklearn2011}
Pedregosa, F. et al. (2011).
Scikit-learn: Machine Learning in Python.
\textit{Journal of Machine Learning Research}, 12, 2825--2830.

\bibitem{pandas2020}
McKinney, W. et al. (2020).
\textit{pandas: a foundational Python library for data analysis and statistics}.
Python for High Performance and Scientific Computing, 14(9).

\bibitem{goodfellow2016}
Goodfellow, I., Bengio, Y. \& Courville, A. (2016).
\textit{Deep Learning}. MIT Press.

\bibitem{hastie2009}
Hastie, T., Tibshirani, R. \& Friedman, J. (2009).
\textit{The Elements of Statistical Learning} (2nd ed.).
Springer.

\bibitem{fielding2000}
Fielding, R.~T. (2000).
\textit{Architectural Styles and the Design of Network-based Software Architectures}.
PhD Thesis, University of California, Irvine.

\bibitem{martin2003}
Martin, R.~C. (2003).
\textit{Agile Software Development: Principles, Patterns, and Practices}.
Prentice Hall.

\end{thebibliography}

\end{document}